\chapter{来源与史料说明及读法索引}
\label{app:sources}

\section{怎样使用本卷材料}

本卷以《宋史》《辽史》《金史》《元史》及相应纪传、志、表建立可复核的年序入口；它们不是任何政权或社会的透明自述。北宋年月日主要回查《续资治通鉴长编》，建炎至绍兴间回查《建炎以来系年要录》，制度条文回查《宋会要辑稿》；三者分别保留编年、政事线索或后出门类，均不能把诏令直接等同于地方执行。跨政权事件还须同对方史料、碑刻、文书和考古材料并读。

\section{基础书目与读法}
\begin{description}
\item[《宋史》] 中华书局点校本。本卷用其本纪、食货志、兵志、地理志和列传建立宋朝政局、制度和宋辽夏金元关系的索引；元末仓促修史、战果数字和人物褒贬使其不能单独承担地方执行、精确军数或对方动机。
\item[《辽史》] 中华书局点校本。本卷用本纪、营卫志、百官志与列传追索辽的帝系、都城、军政和跨境关系；契丹本方连续文书有限，故不能把其制度分类当成辽境社会的全貌。
\item[《金史》] 中华书局点校本。本卷用本纪、食货志、兵志、交聘表与列传编排女真兴起、宋金关系和蒙古战争；金实录及地方档案多佚，华北基层、迁徙者和被征服社会须另以碑志、城址与对方材料限定。
\item[《元史》] 中华书局点校本。本卷用本纪、百官志、食货志、河渠志、兵志和列传建立元朝政治、行省、赋役和河工入口；明初编纂的纪年、姓名与区域效果需与条格、文书、碑刻和多语材料互校。
\item[《续资治通鉴长编》] 中华书局点校本及校勘成果。用于960--1127年宋廷日常、诏令和争议的密集年序；李焘的中枢取舍、传本残缺与引文层次意味着它不能替辽、西夏或民间社会发言。
\item[《建炎以来系年要录》] 中华书局点校本。用于1127--1162年南渡、和战、行在与财政动员的次序；它止于绍兴末，且宋方战报、和议语言和人物判断必须同金方及地方材料并列。
\item[《宋会要辑稿》] 中华书局影印／整理本。用于职官、食货、兵、蕃夷等制度条目的提出、重申与术语追索；徐松辑佚本的错简、缺文和后设门类使其不能证明全国实际执行。
\item[《元典章》《通制条格》] 以可靠整理本为准。用于元代法令、机构和规范语言；条格说明应然秩序，不能独自证明行省、路府或不同身份群体的实际处境。
\item[《三朝北盟会编》与宋人见闻材料] 用于1115年以后宋金冲突、使节与俘虏经验的证据线索；汇编、笔记的追责和传闻成分使日期、军数与动机必须回查编年和对方材料。
\item[《元朝秘史》、波斯文史籍与《高丽史》] 用于蒙古联盟、跨区域征服和高丽关系的并行观察；译名、成书层次、宫廷立场和地名对应都须逐项说明，不能用汉文正史裁定一切。
\item[黑水城西夏文书、契丹／女真／蒙古及汉文碑刻] 用于法律、行政、语言、迁徙、宗教与地方活动的局部证据。须列明文书编号、出土／旧藏、释读版本和定年；集中出土和残片状态限制全国性结论。
\item[宋元契约、港口／沉船、城址与水利考古] 用于土地、市场、海运、生产与河工问题。契约、窖藏、遗址分期和沉船货物各自只能说明特定地点与时间层，不能替代全国统计。
\item[钱穆《国史大纲》《历代政治得失》、吕思勉通史及断代研究] 作为问题意识与制度得失的现代阅读入口，不充当逐年事实、精确数字或人物动机的唯一证据。
\item[《古文观止》] 仅作古文阅读入口；本卷若引选文，须回到可靠原典、版本和成文语境，不能把选本注释当作原始史料。
\end{description}

\section{二十四史的书系定位}

二十四史为《史记》《汉书》《后汉书》《三国志》《晋书》《宋书》《南齐书》《梁书》《陈书》《魏书》《北齐书》《周书》《隋书》《南史》《北史》《旧唐书》《新唐书》《旧五代史》《新五代史》《宋史》《辽史》《金史》《元史》《明史》。本卷直接重点使用其中《宋史》《辽史》《金史》《元史》；其余各书的卷次、版本与适用范围由书系索引另行说明，不能以“二十四史”这一总称抹去材料间的时代、立场和保存差异。

\section{从材料回到叙事}

材料的用途只有回到具体场景才有意义。可由《宋史》《续资治通鉴长编》和《宋会要辑稿》进入\eventref{NS-0960-KAIFENG-CAPITAL}{开封中枢的承接}；由《辽史》、宋人使辽记录与碑志进入\eventref{LIAO-0938-EVENT-031}{幽州的南京经营}；由黑水城文书及多语材料进入\eventref{XIA-0982-EVENT-049}{党项起兵后的河西问题}；由《金史》与宋方材料进入\eventref{JIN-1115-EVENT-079}{金的建国}；由《建炎以来系年要录》进入\eventref{SS-1127-EVENT-111}{南宋行在的重建}；由《元朝秘史》《元史》、波斯文和高丽材料进入\eventref{YUAN-1206-EVENT-151}{草原联盟的重排}与\eventref{YUAN-1271-EVENT-173}{元朝定国号}。更多按政权、问题和图件编排的入口见\hyperref[app:index]{“宋辽夏金元卷读者索引”}。
